\documentclass[10pt,a4paper]{article}

% Paquetes
\usepackage[utf8]{inputenc}
\usepackage[T1]{fontenc}
\usepackage[spanish]{babel}
\usepackage{geometry}
\usepackage{enumitem}
\usepackage{titlesec}
\usepackage{hyperref}

% Diseño de página
\geometry{left=2cm, right=2cm, top=2cm, bottom=2cm}
\pagestyle{empty}

% Formato de secciones - estilo minimal
\titleformat{\section}{\large\bfseries}{}{0em}{}
\titlespacing*{\section}{0pt}{10pt}{4pt}

% Configuración de hipervínculos - sin colores
\hypersetup{
    colorlinks=false,
    pdfborder={0 0 0}
}

% Comandos personalizados
\newcommand{\cvitem}[2]{
    \textbf{#1} \hfill \textit{#2}
}

\newcommand{\cvsubitem}[1]{
    #1
}

\begin{document}

% Encabezado
\begin{center}
    {\LARGE\bfseries Facundo Torraca}

    \vspace{0.2cm}

    facundotorraca@gmail.com $\cdot$ linkedin.com/in/facundotorraca $\cdot$ github.com/facundotorraca
\end{center}

\vspace{0.3cm}

% Resumen Profesional
\section*{Resumen}
Líder de equipo de ingeniería de software con más de 4 años de experiencia especializado en sistemas de resolución de identidades, pipelines de datos y arquitecturas distribuidas. Historial comprobado liderando equipos multifuncionales y diseñando sistemas de alta disponibilidad que aseguran la continuidad del negocio durante transiciones tecnológicas. Experto en ingeniería de datos con dominio de Snowflake, Python, Golang y Kubernetes.

% Experiencia Laboral
\section*{Experiencia}

\cvitem{Líder de Equipo de Ingeniería de Software}{Abr 2024 -- Actualidad}\\
\cvsubitem{Innovid}
\begin{itemize}[leftmargin=*, noitemsep, topsep=2pt, itemsep=2pt]
    \item Lidero un equipo de 3 ingenieros supervisando la resolución de identidades en todos los servicios
    \item Diseño e implemento sistemas de datos de alta disponibilidad procesando millones de señales diarias
    \item Conduzco decisiones arquitectónicas para infraestructura de identidades de publicidad multicanal
    \item Mentoricé miembros del equipo en ingeniería de datos y diseño de sistemas distribuidos
\end{itemize}

\vspace{0.2cm}

\cvitem{Ingeniero de Software Semi-Senior}{May 2022 -- Abr 2024}\\
\cvsubitem{Innovid}
\begin{itemize}[leftmargin=*, noitemsep, topsep=2pt, itemsep=2pt]
    \item Diseñé solución interna de resolución de identidades mejorando precisión en segmentación de anuncios
    \item Aseguré continuidad del negocio implementando estrategias compatibles con privacidad
    \item Construí pipelines de datos integrando con proveedores externos usando Snowflake, DBT y Argo
    \item Lideré desarrollo de servicio de validación de identificadores usando Python
    \item Colaboré con equipos de producto y analítica para definir métricas y KPIs de identidad
\end{itemize}

\vspace{0.2cm}

\cvitem{Desarrollador Full Stack}{Nov 2020 -- Abr 2022}\\
\cvsubitem{Okapii}
\begin{itemize}[leftmargin=*, noitemsep, topsep=2pt, itemsep=2pt]
    \item Desarrollé aplicaciones empresariales personalizadas incluyendo sistemas contables y de inventario
    \item Construí soluciones end-to-end desde requisitos hasta despliegue
    \item Entregué soluciones de software optimizando operaciones comerciales para múltiples clientes
\end{itemize}

% Educación
\section*{Educación}

\cvitem{Ingeniería en Informática (Grado combinado Licenciatura y Maestría)}{2017 -- 2023}\\
\cvsubitem{Universidad de Buenos Aires, Facultad de Ingeniería}

% Habilidades Técnicas
\section*{Habilidades Técnicas}

\textbf{Lenguajes:} Python, Golang, SQL, JavaScript, TypeScript

\textbf{Ingeniería de Datos:} Snowflake, DBT, Argo Workflows, pipelines ETL/ELT

\textbf{Infraestructura:} Kubernetes, Docker, microservicios, CI/CD

\textbf{Dominios:} Resolución de identidades, AdTech, sistemas compatibles con privacidad, sistemas distribuidos

% Proyectos
\section*{Proyectos}

\cvitem{Ticken -- Sistema de Ticketing con Blockchain}{2023}\\
\cvsubitem{github.com/ticken-ts}
\begin{itemize}[leftmargin=*, noitemsep, topsep=2pt, itemsep=2pt]
    \item Diseñé plataforma distribuida de venta de entradas con blockchains públicas y privadas
    \item Construí arquitectura de microservicios en Golang con infraestructura definida en Kubernetes
    \item Proyecto de tesis universitaria demostrando experiencia en sistemas distribuidos
\end{itemize}

% Idiomas
\section*{Idiomas}

Español (Nativo), Inglés (Dominio profesional), Portugués (Lectura y escritura)

\end{document}
